\chapter{Biotechnology News: A New Approach to Treating Depression}

The article is called "A new 'magic mushroom' drug could treat depression without psychedelic hallucinations". It was published by the American Chemical Society on March 8, 2026. This topic is part of biotechnology and medicine. It caught my attention because mental health is a big problem today, even for people working in demanding fields like Information Technology. 
Stress and burnout are common, so new medical solutions are useful.

Many scientists started to look at special mushrooms to help with mental health. People usually call them "magic mushrooms". 
These mushrooms contain a natural chemical called psilocybin. When someone eats the mushroom, their body changes psilocybin into another active chemical called psilocin. 

This chemical strongly affects the brain. It can change how a person feels, sees, and thinks. For a long time, doctors have known that this chemical might help people who suffer from severe sadness (aka clinical depression). 
So they tought that it might also help with other mental health problems like severe anxiety and addiction.

However (there is always a catch), there is a big problem with using these raw mushrooms as medicine. 
When a patient takes them, they experience heavy hallucinations. This means they see and hear things that are not really there. This is also called a "trip".

Many patients are afraid of this side effect. Also, it is very hard for doctors to monitor a patient who is having strong hallucinations. 
Because of this, it is difficult to use this treatment safely in normal hospitals or at home.

Recently, researchers found a smart way to solve this problem using biotechnology. They decided to change the chemical structure of psilocin. 
They used chemistry to create modified versions of the molecule. The main goal was to keep the good medical effects but remove the bad hallucinations.

At the end of the story, they found one very good candidate. They named it "compound 4e". This new molecule is special. When a patient takes it, the body absorbs it very slowly and steadily. Because it is released slowly, it does not shock the brain. It still goes to the brain and connects to the serotonin receptors. 
Compound 4e connects to the exact same spots as the original mushroom chemical.

To prove this works, scientists tested the new compound on mice. They gave some mice the original mushroom chemical, and they gave other mice the new compound 4e. They measured how much of the chemical reached the brain over two days.

To check if the mice were hallucinating, scientists looked at how many times the mice twitched their heads. Head twitching is a known sign of a hallucination in rodents. The mice that took compound 4e had very few head twitches. But the medicine was still working perfectly in their brain cells.

\vspace{0.5cm}
\subsection*{References}
American Chemical Society. (2026, March 8). A new ``magic mushroom'' drug could treat depression without psychedelic hallucinations. \textit{ScienceDaily}. Retrieved March 22, 2026, from https://www.sciencedaily.com/releases/2026/03/260307213232.htm